\section{Төслийн тойм}
\subsection{Шийдэх асуудал ба хамрах хүрээ}
Төсөл нь \textit{Operating System Concepts} сурах бичгийн 8-р бүлэгт авч үзсэн~\cite{osc} олон төрөлт нөөцийн хуваарилалтыг загварчилна. Процесс бүр ажлаа дуусгах хүртэл төрөл тус бүрээс дээд тал нь хэдэн нэгж хэрэглэхээ урьдчилан мэдэгдэнэ. Нөөцийн нэгжүүд тухайн төрөл дотроо ижил. Процесс дуусахад эзэмшиж байсан нэгжүүд нь бусад процесст ашиглагдах боломжтой болно.

БД1-д Safety шалгалт, нөөц хүсэх ба чөлөөлөх үйлдэл, терминалын интерфейс, графын алхамчилсан тоглуулалтыг хэрэгжүүлсэн. БД2-ын зорилго нь нэг өгөгдлийг CLI болон вэбээр туршихад ижил шийдвэр гарч, хэрэглэгч тэр шийдвэрийн шалтгааныг дэлгэцийн тооноос өөрөө тайлбарлаж чаддаг болоход оршино.

Энэ нь хичээлийн дадлагад хоёр ач холбогдолтой. Оюутан хүсэлтийг олгох эсэхийг эхлээд гараар бодож, дараа нь вектор бүрийн өөрчлөлтийг программын хариутай тулгаж чадна. Мөн ижил хүсэлтийг хамгаалалттай болон хамгаалалтгүй горимоор туршихад гардаг ялгааг ажигласнаар сул нөөц хүрэлцэх нь дангаараа аюулгүй байдлыг батлахгүйг ойлгоно.

\subsection{Номын жишээний эх өгөгдөл}
Таван процесс (P0–P4), гурван төрөл (R0, R1, R2) байна. Хүснэгтийн вектор бүр R0, R1, R2 гэсэн тогтмол дараалалтай. \textit{Max} нь процессын мэдэгдсэн дээд хэрэгцээ, \textit{Allocation} нь одоо эзэмшиж буй хэмжээ. Үлдсэн хэрэгцээг нөөцийн төрөл бүрээр хасаж олно.
\[
 \mathrm{Need}_{ij}=\mathrm{Max}_{ij}-\mathrm{Allocation}_{ij}.
\]
\tbl{Номын жишээний анхны хуваарилалт}{tab:init}{
\begin{tabularx}{\linewidth}{YCCC}\toprule
\hd Процесс & \textbf{Max} & \textbf{Allocation} & \textbf{Need}\\\midrule
P0 & $(7,5,3)$ & $(0,1,0)$ & $(7,4,3)$\\
P1 & $(3,2,2)$ & $(2,0,0)$ & $(1,2,2)$\\
P2 & $(9,0,2)$ & $(3,0,2)$ & $(6,0,0)$\\
P3 & $(2,2,2)$ & $(2,1,1)$ & $(0,1,1)$\\
P4 & $(4,3,3)$ & $(0,0,2)$ & $(4,3,1)$\\\bottomrule
\end{tabularx}}

\lead{Эхний төлвийг шалгах нь.} CLI-д өгдөг $(10,5,7)$ нь нийт нөөц. Хүснэгт~\ref{tab:init}-ийн Allocation баганыг төрөл тус бүрээр нэмбэл $(7,2,5)$ гарна. Иймээс сул нэгж буюу $\mathrm{Available}=(10,5,7)-(7,2,5)=(3,3,2)$. Вэб хүснэгтийн Available нүдэнд энэ үлдэгдлийг шууд оруулна. Жишээлбэл, P1-ийн эхний хэрэгцээ R0 дээр $3-2=1$ байна.
\newpage
\section{Эцсийн архитектур}
\subsection{Нэг цөмийг ашигладаг хоёр интерфейс}
Терминалд \code{RQ 1 1 0 2} гэж бичих, браузерт P1-ийг сонгоод хүсэлтийн гурван нүдэнд 1, 0, 2 гэж оруулах нь цөмд ижил үйлдэл хүргэнэ. CLI текстийг процессын дугаар ба тоон вектор болгоно. Вэб интерфейс эдгээрийг JSON хүсэлтээр илгээж, хариуд ирсэн төлвөөр хүснэгт, граф, аюулгүй дарааллыг дахин зурна.

\tbl{Өгөгдөл дамжих зам ба компонентын үүрэг}{tab:arch}{
\begin{tabularx}{\linewidth}{@{}L{32mm}YY@{}}\toprule
\hd Компонент & \textbf{Хүлээж авдаг өгөгдөл} & \textbf{Гүйцэтгэх ажил}\\\midrule
\code{cli.py} & Нийт нөөц, Max ба Allocation файл, командын мөрүүд. & Available-ийг тооцож цөм үүсгэнэ. Команд бүрийн шийдвэрийг хэвлэнэ.\\\midrule
\code{server.py} & Вэбийн хүсэлт, өгөгдөл ба горим. & Оролтыг шалгаад серверийн түгжээ дотор төлөв шинэчилж, үр дүнг буцаана.\\\midrule
\code{core.py} & Батлагдсан төлөв, процессын дугаар, хүсэх буюу буцаах хэмжээ. & Need, аюулгүй дараалал, алхам бүрийн Work-ийг олно. Аюултай түр олголтыг буцаана.\\\midrule
\code{index.html}, \code{graph.js} & Засварлаж буй хүснэгт, серверээс ирсэн төлөв ба алхмууд. & Хүснэгт, хүсэлт, хариу, түүх болон inline SVG графыг харуулна.\\\midrule
\code{scenarios.json} & 3×2 ба 10×5 жишээний матриц. & Вэб ба автомат тест ижил өгөгдлөөр ижил нөхцөлийг давтана.\\\bottomrule
\end{tabularx}}

\subsection{Засварлаж буй ба батлагдсан төлөв}
Вэб хуудсан дээр Available-ийн R0 нүдийг 3-аас 4 болгоход тоо эхлээд зөвхөн хүснэгтийн засварын хувилбарт орно. Тэр агшинд серверт $(3,3,2)$ үлдсэн хэвээр. Ийм хоёр өөр утга харагдаж болох учраас UI анхааруулга үзүүлж Request, Release болон горимын чагтыг түр түгжинэ. «Шалгах» амжилттай болбол серверийн хариу $(4,3,2)$ болж үйлдлүүд дахин нээгдэнэ.

Номын жишээг ачаалах товч, бэлэн жишээ сонгох цэс хоёр \code{loadPreset()} функцийг ашиглана. Тус функц жишээний өгөгдөлтэй хамт \code{enforce} горимыг дамжуулна. Ингэснээр «Номын жишээ» гэсэн сонголтод Safety шалгалт асаж, хүсэлтийн талбарууд тэглэгдэнэ. Banker-гүй туршилтыг үргэлжлүүлэх бол тухайн жишээг тусад нь сонгоно.
\lead{P1-ийн хүсэлтээр мөрдөх нь.} Вэбээс \code{/api/request} руу процесс 1, хүсэлт $(1,0,2)$ ирэхэд серверийн \code{request()} цөмийн шийдвэрийг авна. Хариуны \code{request_result} хэсэгт «олгосон» эсэх, \code{state} хэсэгт олголтын дараах хуваарилалт ба аюулгүй дараалал ирнэ. Дэлгэцийн мэдэгдэл эхнийхийг, хүснэгт ба граф хоёр дахь хэсгийг уншдаг. Иймээс зөвшөөрөгдөөгүй хүсэлтэд улаан мэдэгдэл гарсан ч аюулгүй төлвийн тэмдэг ногоон хэвээр байж болно.
\newpage
